% Esquema del pipeline GNN-HVA (figura TikZ).
% Requiere en el preambulo: \usepackage{tikz}, \usepackage{xcolor},
% \usetikzlibrary{arrows.meta, positioning, shapes.geometric, fit, calc},
% y % Paleta de colores estandarizada de la tesis GNN-HVA.
% Fuente de verdad unica: toda figura (TikZ, pgfplots) y toda tabla que use
% color debe importar este archivo con % Paleta de colores estandarizada de la tesis GNN-HVA.
% Fuente de verdad unica: toda figura (TikZ, pgfplots) y toda tabla que use
% color debe importar este archivo con % Paleta de colores estandarizada de la tesis GNN-HVA.
% Fuente de verdad unica: toda figura (TikZ, pgfplots) y toda tabla que use
% color debe importar este archivo con \input{tesis-figures/palette} y usar
% los nombres semanticos, nunca HEX sueltos.
%
% Requiere \usepackage{xcolor} en el preambulo del documento.
%
% Las cuatro etapas del pipeline forman una progresion neutro -> destacado:
%   1. Definicion del problema   gris          #8A8D91
%   2. Datos de referencia       amarillo suave #F2C94C
%   3. Optimizacion VQE          naranja        #E8833A
%   4. Prediccion y evaluacion   violeta        #7B54B8
%
% Auxiliares para resultados (buen/mal resultado), accesibles en escala de grises.

% Etapas del pipeline
\definecolor{stageProblem}{HTML}{8A8D91}  % gris    -- etapa 1
\definecolor{stageData}{HTML}{F2C94C}     % amarillo -- etapa 2
\definecolor{stageVQE}{HTML}{E8833A}      % naranja -- etapa 3
\definecolor{stagePred}{HTML}{7B54B8}     % violeta -- etapa 4

% Auxiliares de resultado
\definecolor{resultGood}{HTML}{2E86AB}    % azul    -- configuracion que aprueba
\definecolor{resultBad}{HTML}{C0392B}     % rojo    -- configuracion que falla

% Tintes claros para rellenos de fondo (borde en el color pleno, relleno tenue)
\definecolor{stageProblemBg}{HTML}{E8E9EA}
\definecolor{stageDataBg}{HTML}{FBEFC7}
\definecolor{stageVQEBg}{HTML}{F9DDC7}
\definecolor{stagePredBg}{HTML}{E4DAF2}
 y usar
% los nombres semanticos, nunca HEX sueltos.
%
% Requiere \usepackage{xcolor} en el preambulo del documento.
%
% Las cuatro etapas del pipeline forman una progresion neutro -> destacado:
%   1. Definicion del problema   gris          #8A8D91
%   2. Datos de referencia       amarillo suave #F2C94C
%   3. Optimizacion VQE          naranja        #E8833A
%   4. Prediccion y evaluacion   violeta        #7B54B8
%
% Auxiliares para resultados (buen/mal resultado), accesibles en escala de grises.

% Etapas del pipeline
\definecolor{stageProblem}{HTML}{8A8D91}  % gris    -- etapa 1
\definecolor{stageData}{HTML}{F2C94C}     % amarillo -- etapa 2
\definecolor{stageVQE}{HTML}{E8833A}      % naranja -- etapa 3
\definecolor{stagePred}{HTML}{7B54B8}     % violeta -- etapa 4

% Auxiliares de resultado
\definecolor{resultGood}{HTML}{2E86AB}    % azul    -- configuracion que aprueba
\definecolor{resultBad}{HTML}{C0392B}     % rojo    -- configuracion que falla

% Tintes claros para rellenos de fondo (borde en el color pleno, relleno tenue)
\definecolor{stageProblemBg}{HTML}{E8E9EA}
\definecolor{stageDataBg}{HTML}{FBEFC7}
\definecolor{stageVQEBg}{HTML}{F9DDC7}
\definecolor{stagePredBg}{HTML}{E4DAF2}
 y usar
% los nombres semanticos, nunca HEX sueltos.
%
% Requiere \usepackage{xcolor} en el preambulo del documento.
%
% Las cuatro etapas del pipeline forman una progresion neutro -> destacado:
%   1. Definicion del problema   gris          #8A8D91
%   2. Datos de referencia       amarillo suave #F2C94C
%   3. Optimizacion VQE          naranja        #E8833A
%   4. Prediccion y evaluacion   violeta        #7B54B8
%
% Auxiliares para resultados (buen/mal resultado), accesibles en escala de grises.

% Etapas del pipeline
\definecolor{stageProblem}{HTML}{8A8D91}  % gris    -- etapa 1
\definecolor{stageData}{HTML}{F2C94C}     % amarillo -- etapa 2
\definecolor{stageVQE}{HTML}{E8833A}      % naranja -- etapa 3
\definecolor{stagePred}{HTML}{7B54B8}     % violeta -- etapa 4

% Auxiliares de resultado
\definecolor{resultGood}{HTML}{2E86AB}    % azul    -- configuracion que aprueba
\definecolor{resultBad}{HTML}{C0392B}     % rojo    -- configuracion que falla

% Tintes claros para rellenos de fondo (borde en el color pleno, relleno tenue)
\definecolor{stageProblemBg}{HTML}{E8E9EA}
\definecolor{stageDataBg}{HTML}{FBEFC7}
\definecolor{stageVQEBg}{HTML}{F9DDC7}
\definecolor{stagePredBg}{HTML}{E4DAF2}
.
%
% Fiel a la tesis (tres fases + entrada):
%   Entrada (gris)  Fase 1 datos de referencia (amarillo)
%   Fase 2 VQE (naranja)  Fase 3 prediccion y evaluacion (violeta)
% Recuadro de entrenamiento de la GNN debajo: se entrena con el grafo del
% Hamiltoniano (entrada) y las etiquetas theta* = parametros optimos del
% circuito HVA (salida de Fase 2). Ya entrenada, predice en Fase 3, donde el
% mismo circuito HVA se reutiliza para evaluar la energia con theta_pred.
% El rotulo de fase va ARRIBA del recuadro, en dos lineas.

\begin{tikzpicture}[
    font=\small,
    stage/.style={
        rectangle, rounded corners=3pt, draw=black, line width=1pt,
        fill=#1, minimum width=3.6cm, minimum height=4.3cm, align=center,
        inner sep=6pt,
    },
    gnnbox/.style={
        rectangle, rounded corners=3pt, draw=black, line width=1pt,
        fill=stagePredBg, minimum width=6.4cm, minimum height=3.1cm,
        align=center, inner sep=6pt,
    },
    phaselabel/.style={font=\footnotesize\bfseries, align=center, text=black},
    subtag/.style={font=\scriptsize, text=black!70, align=center},
    outtag/.style={font=\scriptsize\ttfamily, text=black!55, align=center},
    flow/.style={-{Stealth[length=3.2mm]}, line width=1.2pt, black!80},
    trainflow/.style={-{Stealth[length=3mm]}, line width=1.1pt, dashed,
        stagePred!55!black},
    flowlabel/.style={font=\scriptsize\itshape, text=black!75, align=center,
        fill=white, inner sep=2pt},
    icon/.style={line width=0.9pt},
    node distance=1.55cm,
]

% ===== ENTRADA: Definicion del problema (gris) =====
\node[stage=stageProblemBg] (s1) {};
\node[phaselabel] at ([yshift=0.62cm]s1.north) {Entrada};
\node[phaselabel, font=\scriptsize\bfseries] at ([yshift=0.28cm]s1.north)
    {Definici\'on del problema};
\begin{scope}[shift={($(s1.center)+(0,0.9cm)$)}]
    \foreach \x/\y in {-0.75/0.2, -0.05/0.45, 0.7/0.25, -0.45/-0.4, 0.4/-0.45} {
        \node[circle, fill=stageProblem, draw=black, line width=0.5pt,
              minimum size=6pt, inner sep=0pt] at (\x,\y) {};
    }
    \draw[icon, black!70] (-0.75,0.2)--(-0.05,0.45)--(0.7,0.25);
    \draw[icon, black!70] (-0.75,0.2)--(-0.45,-0.4)--(0.4,-0.45)--(0.7,0.25);
    \draw[icon, black!70] (-0.05,0.45)--(-0.45,-0.4);
\end{scope}
\node[subtag] at ($(s1.center)+(0,-0.5cm)$)
    {Hamiltoniano $H(h)$\\sobre una topolog\'ia\\(red de espines)};
\node[outtag] at ($(s1.center)+(0,-1.6cm)$) {$-J\!\sum\! ZZ - h\!\sum\! X$};

% ===== FASE 1: Datos de referencia (amarillo) =====
\node[stage=stageDataBg, right=of s1] (s2) {};
\node[phaselabel] at ([yshift=0.62cm]s2.north) {Fase 1};
\node[phaselabel, font=\scriptsize\bfseries] at ([yshift=0.28cm]s2.north)
    {Datos de referencia};
\begin{scope}[shift={($(s2.center)+(0,0.9cm)$)}]
    \draw[icon, black!70, -{Stealth[length=1.6mm]}] (-0.8,-0.45)--(0.85,-0.45);
    \draw[icon, black!70, -{Stealth[length=1.6mm]}] (-0.8,-0.45)--(-0.8,0.55);
    \node[font=\tiny, text=black!60] at (0.8,-0.62) {$h$};
    \node[font=\tiny, text=black!60] at (-0.97,0.45) {$E_0$};
    \draw[icon, stageData!55!black, line width=1.3pt]
        (-0.75,0.35) .. controls (-0.2,-0.3) and (0.25,-0.38) .. (0.8,-0.18);
\end{scope}
\node[subtag] at ($(s2.center)+(0,-0.5cm)$)
    {Diagonalizaci\'on exacta\\o DMRG: estado\\fundamental por cada $h$};
\node[outtag] at ($(s2.center)+(0,-1.55cm)$)
    {$E_0,\ \mathrm{gap},$\\$\langle X_i\rangle,\ \langle Z_iZ_{i+1}\rangle$};

% ===== FASE 2: Optimizacion VQE (naranja) =====
\node[stage=stageVQEBg, right=of s2] (s3) {};
\node[phaselabel] at ([yshift=0.62cm]s3.north) {Fase 2};
\node[phaselabel, font=\scriptsize\bfseries] at ([yshift=0.28cm]s3.north)
    {VQE con warm-start};
\begin{scope}[shift={($(s3.center)+(0,0.9cm)$)}]
    \draw[icon, black!75] (-0.9,0.25)--(0.9,0.25);
    \draw[icon, black!75] (-0.9,-0.25)--(0.9,-0.25);
    \node[font=\tiny, text=black!55] at (-1.08,0.25) {$q_0$};
    \node[font=\tiny, text=black!55] at (-1.08,-0.25) {$q_1$};
    \node[draw=black, fill=stageVQE!35, minimum size=0.32cm, inner sep=0pt,
          font=\tiny] at (-0.45,0.25) {$R_x$};
    \node[draw=black, fill=stageVQE!35, minimum size=0.32cm, inner sep=0pt,
          font=\tiny] at (-0.45,-0.25) {$R_x$};
    \filldraw[black] (0.2,0.25) circle (1.6pt);
    \filldraw[black] (0.2,-0.25) circle (1.6pt);
    \draw[icon, black] (0.2,0.25)--(0.2,-0.25);
    \node[draw=black, fill=stageVQE!35, minimum size=0.3cm, inner sep=0pt,
          font=\tiny] at (0.6,0.25) {$R_z$};
    \node[font=\tiny, text=black!55] at (0,-0.62) {circuito HVA};
\end{scope}
\node[subtag] at ($(s3.center)+(0,-0.55cm)$)
    {Se optimizan los \'angulos\\del circuito HVA\\(barrido descendente en $h$)};
\node[outtag] at ($(s3.center)+(0,-1.6cm)$)
    {$\theta^{*}(h)$ (\'angulos HVA)};

% ===== FASE 3: Prediccion y evaluacion (violeta) =====
\node[stage=stagePredBg, right=of s3] (s4) {};
\node[phaselabel] at ([yshift=0.62cm]s4.north) {Fase 3};
\node[phaselabel, font=\scriptsize\bfseries] at ([yshift=0.28cm]s4.north)
    {Predicci\'on y evaluaci\'on};
\begin{scope}[shift={($(s4.center)+(0,0.9cm)$)}]
    \node[circle, draw=black, fill=stagePred!20, minimum size=6pt, inner sep=0pt] (p1) at (-0.9,0.28) {};
    \node[circle, draw=black, fill=stagePred!20, minimum size=6pt, inner sep=0pt] (p2) at (-0.9,-0.28) {};
    \node[draw=black, fill=stagePred!45, rounded corners=2pt,
          minimum width=0.75cm, minimum height=0.5cm, font=\tiny, inner sep=1pt]
          (gnnmini) at (0.0,0.0) {GNN};
    \draw[icon, black!70, -{Stealth[length=1.4mm]}] (p1)--(gnnmini.west);
    \draw[icon, black!70, -{Stealth[length=1.4mm]}] (p2)--(gnnmini.west);
    \draw[icon, black!70, -{Stealth[length=1.8mm]}] (gnnmini.east)--(0.95,0.0);
    \node[font=\tiny, text=stagePred!60!black] at (0.78,0.26) {$\theta_{\mathrm{pred}}$};
\end{scope}
\node[subtag] at ($(s4.center)+(0,-0.55cm)$)
    {La GNN predice $\theta$ en\\una inferencia; el circuito\\HVA eval\'ua la energ\'ia};
\node[outtag] at ($(s4.center)+(0,-1.6cm)$) {$\Delta E/\mathrm{gap} < 5\%$};

% ===== Flechas de flujo (fila superior) =====
\draw[flow] (s1) -- (s2) node[flowlabel, midway, above] {grafo de $H$};
\draw[flow] (s2) -- (s3) node[flowlabel, midway, above] {$E_0(h),$ gap};

% ===== Recuadro: Entrenamiento de la GNN, debajo y centrado =====
\coordinate (midbelow) at ($(s2.south)!0.5!(s3.south)$);
\node[gnnbox, anchor=north] (gnn) at ($(midbelow)+(0,-3.0cm)$) {};
\node[phaselabel] at ([yshift=0.42cm]gnn.north)
    {Entrenamiento de la GNN (MPNN)};
% arquitectura MPNN a la izquierda
\begin{scope}[shift={($(gnn.center)+(-2.0cm,0.25cm)$)}]
    \node[circle, draw=black, fill=stagePred!25, minimum size=6pt, inner sep=0pt] (n1) at (-0.5,0.35) {};
    \node[circle, draw=black, fill=stagePred!25, minimum size=6pt, inner sep=0pt] (n2) at (-0.5,-0.35) {};
    \node[circle, draw=black, fill=stagePred!45, minimum size=7pt, inner sep=0pt] (n3) at (0.15,0.0) {};
    \draw[icon, black!70] (n1)--(n3); \draw[icon, black!70] (n2)--(n3);
\end{scope}
\node[font=\tiny, text=black!60, align=center]
    at ($(gnn.center)+(-1.85cm,-0.72cm)$) {3$\times$ GINConv + pooling};
% objetivo de entrenamiento, a la derecha
\node[align=center, font=\scriptsize, text=black!75]
    at ($(gnn.center)+(1.35cm,0.35cm)$)
    {aprende $h \mapsto \theta$\\minimizando MSE\\$\theta_{\mathrm{pred}} \approx \theta^{*}$};
\node[subtag, font=\scriptsize, text=black!60]
    at ($(gnn.center)+(1.35cm,-0.75cm)$)
    {entrada: grafo de $H$\\(nodos $=h_i$, aristas $=J_{ij}$)};

% ===== Flechas de datos de entrenamiento (ortogonales) =====
% grafo del Hamiltoniano (desde Fase 1) -> entra por la izquierda del recuadro
\draw[trainflow] (s2.south) |- ($(gnn.west)+(0,0.15cm)$)
    node[flowlabel, pos=0.72] {grafo + observables};
% etiquetas theta* = angulos del circuito HVA (desde Fase 2) -> entra por arriba
\draw[trainflow] (s3.south) |- ($(gnn.north)+(1.1cm,0)$)
    node[flowlabel, pos=0.72] {$\theta^{*}$ (etiquetas)};

% ===== GNN entrenada -> Fase 3 (predice los angulos) =====
\draw[flow, stagePred!55!black]
    (gnn.east) -| (s4.south)
    node[flowlabel, pos=0.28] {GNN entrenada};

\end{tikzpicture}
